% answers_short.tex —— 2011 年·山东卷·理科 · 简短答案速查
% 由 tools/make_short_answers.py 自 answers.tex 的 \gksolhead/\gkanswer 提取，
% 并与 2011/2011_数学真题_山东理_答案要点.md 逐题比对后生成，勿手改（改源 MD 后重跑）。

\begin{gknotice}
\textbf{使用说明}：本节只给最终答案便于快速核对，完整推导见解析版。答案为编写者独立推导，\textbf{非官方标准答案}。
\end{gknotice}

\gkpart{选择题}
\begin{center}\begin{tabular}{ccccccccccccc}
\toprule
题号 & 1 & 2 & 3 & 4 & 5 & 6 & 7 & 8 & 9 & 10 & 11 & 12 \\
\midrule
答案 & A & D & D & D & B & C & B & A & C & B & A & D \\
\bottomrule
\end{tabular}\end{center}

\gkpart{填空题}
\begin{itemize}[leftmargin=2.6em]
  \item[\textbf{13.}] $68$。
  \item[\textbf{14.}] $4$。
  \item[\textbf{15.}] $\dfrac{x}{(2^n-1)x+2^n}$。
  \item[\textbf{16.}] $2$。
\end{itemize}

\gkpart{解答题}
\begin{itemize}[leftmargin=2.6em]
  \item[\textbf{17.}] (1) $2$；(2) $\dfrac{\sqrt{15}}{4}$
  \item[\textbf{18.}] (1) $0.55$；(2) $P(\xi=0)=0.1$，$P(\xi=1)=0.35$，$P(\xi=2)=0.4$，$P(\xi=3)=0.15$，$E\xi=1.6$
  \item[\textbf{19.}] (1) 证明见解析；(2) $60^\circ$（即 $\dfrac{\pi}{3}$）
  \item[\textbf{20.}] (1) $a_n=2\cdot3^{n-1}$；(2) $n$ 为偶数时 $S_n=3^n-1+\dfrac n2\ln3$；$n$ 为奇数时 $S_n=3^n-1-\ln2-\dfrac{n-1}{2}\ln3$
  \item[\textbf{21.}] (1) $y=\dfrac{160\pi}{r}+4\pi(c-2)r^2$，定义域 $0<r\leqslant2$；(2) $3<c\leqslant4.5$ 时 $r=2$；$c>4.5$ 时 $r=\sqrt[3]{\dfrac{20}{c-2}}$
  \item[\textbf{22.}] (1) 证明见解析，$x_1^2+x_2^2=3$，$y_1^2+y_2^2=2$；(2) 最大值 $\dfrac52$；(3) 不存在
\end{itemize}
